\section{Dimensionamento dos componentes}

A primeira parte do projeto consiste na dimensionamento das grandezas dos 4 componentes princiais do conversor. Todos devem ser obtidos por meio das deduções matemáticas a seguir visando atender aos critérios de projeto apresentados na tabela \ref{tab:especificacoes-buck}.

\begin{table}[ht]
    \centering
    \caption{Especificações do conversor Buck.}
    \label{tab:especificacoes-buck}
    \begin{tabular}{lcc}
        \toprule
        \textbf{Parâmetro} & \textbf{Símbolo} & \textbf{Valor} \\
        \midrule

        Tensão de entrada
            & $V_{\mathrm{IN}}$
            & \SI{24}{\volt} \\

        Tensão de saída
            & $V_{\mathrm{OUT}}$
            & \SI{12}{\volt} \\

        Potência de saída
            & $P_{\mathrm{OUT}}$
            & \SI{40}{\watt} \\

        Rendimento
            & $\eta$
            & 95--100\,\% \\

        Frequência de chaveamento
            & $f_{\mathrm{S}}$
            & \SI{50}{\kilo\hertz} \\

        Ondulação de corrente no indutor
            & $\Delta I_{\mathrm{L}}$
            & 5\,\% \\

        Ondulação de tensão no capacitor
            & $\Delta V_{\mathrm{C}}$
            & 1\,\% \\

        \bottomrule
    \end{tabular}

\end{table}

Os primeiros parametros a serem determinados são o Duty Cycle e a corrente de saída, pois dependem apenas das tensões e da potência de projeto. As equações \ref{eq:duty-cycle-proj} e \ref{eq:corrente-saida} sintetizam os cálculos.

\begin{equation}
    D = \frac{V_{\mathrm{OUT}}}{V_{\mathrm{IN}}}
    = \frac{12}{24}
    = 0{,}5
    \quad\Rightarrow\quad
    \boxed{D = 50\,\%}
    \label{eq:duty-cycle-proj}
\end{equation}

\begin{equation}
    I_{\mathrm{OUT}} =
    \frac{P_{\mathrm{OUT}}}{V_{\mathrm{OUT}}}=
    \frac{40}{12}
    \approx 3{,}33\,\mathrm{A}
    \quad\Rightarrow\quad
    \boxed{I_{\mathrm{OUT}} \approx 3{,}33\,\mathrm{A}}
    \label{eq:corrente-saida}
\end{equation}
  
Com a corrente de saída calculada, pode-se encontrar a ondulação de corrente máxima aceitável por critério de projeto. Ela corresponde a diferença entre a corrente de pico e vale da saída, como mostra a figura \ref{fig:ondulacao-corrente-indutor}. Ambos os valores são calculados nas equações \ref{eq:ondulacao-corrente}, \ref{eq:corrente-pico} e \ref{eq:corrente-vale}.

\begin{figure}[h]
    \centering
    \caption{Ondulação da corrente no indutor de um conversor Buck.}

    \begin{tikzpicture}
        \begin{axis}[
            width=0.9\textwidth,
            height=7cm,
            axis lines=middle,
            xlabel={$t$},
            ylabel={$i_L(t)$},
            xmin=0,
            xmax=4,
            ymin=0,
            ymax=7,
            xtick=\empty,
            ytick=\empty,
            axis line style={->},
            xlabel style={anchor=west},
            ylabel style={anchor=south},
            clip=false
        ]

        % Corrente média do indutor
        \addplot[
            dashed,
            domain=0:4,
            samples=2
        ] {4};

        \addplot[
            dashed,
            domain=0:4,
            samples=2
        ] {3.5};

        \addplot[
            dashed,
            domain=0:4,
            samples=2
        ] {3};
        
        % Corrente do indutor
        \addplot[
            thick,
            domain=0:1,
            samples=2
        ] {3 + x};

        \addplot[
            thick,
            domain=1:2,
            samples=2
        ] {5 - x};

        \addplot[
            thick,
            domain=2:3,
            samples=2
        ] {3 + x - 2};

        \addplot[
            thick,
            domain=3:4,
            samples=2
        ] {5 - x + 2};

        % Linha vertical indicando a ondulação
        \draw[<->, thick]
            (axis cs:1,3) -- 
            node[right=4pt] {$\Delta I_L$}
            (axis cs:1,4);

        % Marca do pico
        \node[above] at (axis cs:4,4) {$I_{L,\mathrm{máx}}$};

        % Marca do vale
        \node[below] at (axis cs:4,3) {$I_{L,\mathrm{mín}}$};

        % Indicação da corrente média
        \node[right] at (axis cs:4,3.5) {$I_{L,\mathrm{média}}$};

        \end{axis}
    \end{tikzpicture}
    \label{fig:ondulacao-corrente-indutor}
\end{figure}

\begin{equation}
    \Delta I_L = 0{,}05 \cdot I_{\mathrm{OUT}}
    \quad\Rightarrow\quad
    \Delta I_L = 0{,}05 \cdot 3{,}33
    \approx 0{,}167\,\mathrm{A}
    \quad\Rightarrow\quad
    \boxed{\Delta I_L \approx 0{,}167\,\mathrm{A}}
    \label{eq:ondulacao-corrente}
\end{equation}

\begin{equation}
    I_{L,\mathrm{pico}} =
    I_{\mathrm{OUT}} + \frac{\Delta I_L}{2}
    \quad\Rightarrow\quad
    3{,}33 + \frac{0{,}167}{2}
    \quad\Rightarrow\quad
    \boxed{I_{L,\mathrm{pico}}\approx 3{,}41\,\mathrm{A}}
    \label{eq:corrente-pico}
\end{equation}

\begin{equation}
    I_{L,\mathrm{vale}} =
    I_{\mathrm{OUT}} - \frac{\Delta I_L}{2}
    \quad\Rightarrow\quad
    3{,}33 - \frac{0{,}167}{2}
    \quad\Rightarrow\quad
    \boxed{I_{L,\mathrm{vale}}\approx 3{,}25\,\mathrm{A}}
    \label{eq:corrente-vale}
\end{equation}

A indutância do Indutor pode ser determinada utilizando a equação \ref{eq:indutancia-buck}, todos os valorese necessários para a derteminação da mesma foram calculados nas equações anteriores. Substituindo eles obtemos o resultado na equação \ref{eq:valor-indutancia}.


\begin{equation}
    L =
    \frac{(V_{\mathrm{IN}}-V_{\mathrm{OUT}})D}
    {\Delta I_L f_S}
    \label{eq:indutancia-buck}
\end{equation}

\begin{equation}
    L =
    \frac{(24-12)\cdot0{,}5}
    {0{,}167\cdot50\times10^3}
    \approx 0{,}718\,\mathrm{mH}
    \quad\Rightarrow\quad
    \boxed{L \approx 718\,\mu\mathrm{H}}
    \label{eq:valor-indutancia}
\end{equation}


Assim encerra o processo de dimensionar o indutor, essas sçao as informações dele compiladas:

\begin{table}[ht]
    \centering
    \caption{Parâmetros de projeto do indutor do conversor Buck.}
    \label{tab:parametros-indutor}
    \begin{tabular}{lc}
        \toprule
        \textbf{Parâmetro} & \textbf{Valor} \\
        \midrule
        Tensão de entrada, $V_{\mathrm{IN}}$ 
            & \SI{24}{\volt} \\
        Tensão de saída, $V_{\mathrm{OUT}}$ 
            & \SI{12}{\volt} \\
        Potência de saída, $P_{\mathrm{OUT}}$ 
            & \SI{40}{\watt} \\
        Frequência de chaveamento, $f_{\mathrm{S}}$ 
            & \SI{50}{\kilo\hertz} \\
        Ciclo de trabalho, $D$ 
            & 50\,\% \\
        Corrente média, $I_{\mathrm{OUT}}$ 
            & \SI{3.33}{\ampere} \\
        Ondulação de corrente, $\Delta I_{\mathrm{L}}$ 
            & \SI{0.167}{\ampere} \;(5\,\%) \\
        Corrente de pico, $I_{\mathrm{L,pico}}$ 
            & \SI{3.41}{\ampere} \\
        Corrente de vale, $I_{\mathrm{L,vale}}$ 
            & \SI{3.25}{\ampere} \\
        Indutância calculada, $L$ 
            & \SI{718}{\micro\henry} \\
        Indutância selecionada 
            & \SI{750}{\micro\henry} \\
        \bottomrule
    \end{tabular}
\end{table}

\subsection{Capacitor de saída}

O capacitância de saída do indutor é encontrada por meio da equação \ref{eq:valor-capacitancia}.

\begin{equation}
    \Delta V_{\mathrm{OUT}} = 0{,}01 \cdot V_{\mathrm{OUT}}
    \quad\Rightarrow\quad
    \Delta V_{\mathrm{OUT}} = 0{,}01 \cdot 12 V
    = 0{,}12\,\mathrm{V}
    \quad\Rightarrow\quad
    \boxed{\Delta V_{\mathrm{OUT}} = 0{,}12\,\mathrm{V}}
    \label{eq:ondulacao-corrente}
\end{equation}

\begin{equation}
    C_{\mathrm{OUT}} =
    \frac{\Delta I_L}
    {8\cdot f_s \cdot \Delta V_{OUT}} =
    \frac{0{,}167}
    {8\cdot 50 \times 10^{5} \cdot 0{,}12}
    \approx 3{,}48\,\mu\mathrm{F}
    \quad\Rightarrow\quad
    \boxed{C_{\mathrm{OUT}} \approx 3{,}48\,\mu\mathrm{F}}
    \label{eq:valor-capacitancia}
\end{equation}

\subsection{MOSFET}

O MOSFET deve ser capaz de aguntear a tensão de entrada com uma margem de seguranaça, adotada em 1,5, ou seja a tensão entre o drain e o Source deve ser

\begin{equation}
    V_{DS}=V_{IN} \cdot 1,5 = 24 \cdot 1,5 = 36V
\end{equation}

Já a corrente que o MOSFET deve aguentar é a corrente eficaz de saida

\begin{equation}
    I_{Drms} = I_{OUT}*\sqrt{D}=3,33 \cdot \sqrt{0,5} = 2,35A
\end{equation}

\subsection{Diodo}

Já o diodo deve ser capaz de aguntar uma corrente e uma tensão de polarização reversa similar a de VDS

